% Options for packages loaded elsewhere
\PassOptionsToPackage{unicode}{hyperref}
\PassOptionsToPackage{hyphens}{url}
\documentclass[
  12pt,
  ignorenonframetext,
  aspectratio=169]{beamer}
\newif\ifbibliography
\usepackage{pgfpages}
\setbeamertemplate{caption}[numbered]
\setbeamertemplate{caption label separator}{: }
\setbeamercolor{caption name}{fg=normal text.fg}
\beamertemplatenavigationsymbolsempty
% remove section numbering
\setbeamertemplate{part page}{
  \centering
  \begin{beamercolorbox}[sep=16pt,center]{part title}
    \usebeamerfont{part title}\insertpart\par
  \end{beamercolorbox}
}
\setbeamertemplate{section page}{
  \centering
  \begin{beamercolorbox}[sep=12pt,center]{section title}
    \usebeamerfont{section title}\insertsection\par
  \end{beamercolorbox}
}
\setbeamertemplate{subsection page}{
  \centering
  \begin{beamercolorbox}[sep=8pt,center]{subsection title}
    \usebeamerfont{subsection title}\insertsubsection\par
  \end{beamercolorbox}
}
% Prevent slide breaks in the middle of a paragraph
\widowpenalties 1 10000
\raggedbottom
\AtBeginPart{
  \frame{\partpage}
}
\AtBeginSection{
  \ifbibliography
  \else
    \frame{\sectionpage}
  \fi
}
\AtBeginSubsection{
  \frame{\subsectionpage}
}
\usepackage{iftex}
\ifPDFTeX
  \usepackage[T1]{fontenc}
  \usepackage[utf8]{inputenc}
  \usepackage{textcomp} % provide euro and other symbols
\else % if luatex or xetex
  \usepackage{unicode-math} % this also loads fontspec
  \defaultfontfeatures{Scale=MatchLowercase}
  \defaultfontfeatures[\rmfamily]{Ligatures=TeX,Scale=1}
\fi
\usepackage{lmodern}
\ifPDFTeX\else
  % xetex/luatex font selection
\fi
% Use upquote if available, for straight quotes in verbatim environments
\IfFileExists{upquote.sty}{\usepackage{upquote}}{}
\IfFileExists{microtype.sty}{% use microtype if available
  \usepackage[]{microtype}
  \UseMicrotypeSet[protrusion]{basicmath} % disable protrusion for tt fonts
}{}
\makeatletter
\@ifundefined{KOMAClassName}{% if non-KOMA class
  \IfFileExists{parskip.sty}{%
    \usepackage{parskip}
  }{% else
    \setlength{\parindent}{0pt}
    \setlength{\parskip}{6pt plus 2pt minus 1pt}}
}{% if KOMA class
  \KOMAoptions{parskip=half}}
\makeatother
\setlength{\emergencystretch}{3em} % prevent overfull lines
\providecommand{\tightlist}{%
  \setlength{\itemsep}{0pt}\setlength{\parskip}{0pt}}
\usepackage{lmodern}
\definecolor{ink}{HTML}{1F2933}
\definecolor{teal}{HTML}{0F766E}
\definecolor{orange}{HTML}{D95D39}
\definecolor{muted}{HTML}{667085}
\setbeamercolor{normal text}{fg=ink,bg=white}
\setbeamercolor{structure}{fg=teal}
\setbeamercolor{frametitle}{fg=ink,bg=white}
\setbeamercolor{title}{fg=ink}
\setbeamercolor{subtitle}{fg=muted}
\setbeamercolor{alerted text}{fg=orange}
\setbeamercolor{block title}{fg=white,bg=teal}
\setbeamercolor{block body}{fg=ink,bg=teal!6}
\setbeamerfont{title}{series=\bfseries,size=\LARGE}
\setbeamerfont{frametitle}{series=\bfseries,size=\Large}
\setbeamertemplate{navigation symbols}{}
\setbeamertemplate{headline}{}
\setbeamertemplate{frametitle continuation}{}
\setbeamertemplate{footline}{\hfill\scriptsize\color{muted}\insertframenumber\, / \,\inserttotalframenumber\hspace{0.35cm}\vspace{0.18cm}}
\newcommand{\sourcecite}[1]{\vfill{\tiny\color{muted}#1}}
\usepackage{bookmark}
\IfFileExists{xurl.sty}{\usepackage{xurl}}{} % add URL line breaks if available
\urlstyle{same}
\hypersetup{
  pdftitle={Quando a IA faz a análise, quem responde pelo método?},
  pdfauthor={Manoel Galdino},
  hidelinks,
  pdfcreator={LaTeX via pandoc}}

\title{Quando a IA faz a análise, quem responde pelo método?}
\subtitle{Abstração, dívida epistêmica e defesa metodológica sem cola}
\author{Manoel Galdino}
\date{7 de agosto de 2026}

\begin{document}
\frame{\titlepage}

\begin{frame}{A pergunta que eu não consegui responder}
\protect\phantomsection\label{a-pergunta-que-eu-nuxe3o-consegui-responder}
\begin{center}
\vspace{0.6cm}
{\color{orange}\Huge\textbf{Eu descobri que não sabia defender uma decisão do meu próprio paper.}}

\vspace{0.8cm}
\large Não era um erro de código. Era uma perda de autoria metodológica.
\end{center}
\end{frame}

\begin{frame}{Antes de falar de IA, pense no seu método}
\protect\phantomsection\label{antes-de-falar-de-ia-pense-no-seu-muxe9todo}
\begin{center}
\vspace{0.6cm}
\Large\textbf{Qual foi a última decisão metodológica que você conseguiria reconstruir hoje - sem abrir o código?}

\vspace{0.9cm}
\normalsize O que você teria de explicar?

\vspace{0.25cm}
\begin{tabular}{c@{\qquad}c@{\qquad}c}
  objeto & medida & interpretação \\
  \color{teal}\rule{2.2cm}{1.2pt} & \color{teal}\rule{2.2cm}{1.2pt} & \color{teal}\rule{2.2cm}{1.2pt}
\end{tabular}
\end{center}
\end{frame}

\begin{frame}{A IA é uma nova camada de abstração}
\protect\phantomsection\label{a-ia-uxe9-uma-nova-camada-de-abstrauxe7uxe3o}
\begin{center}
\vspace{0.2cm}
\large
\texttt{linguagem de máquina}

\vspace{0.05cm}
{\color{teal}\large\downarrow}

\vspace{0.05cm}
\texttt{linguagem de programação}

\vspace{0.05cm}
{\color{teal}\large\downarrow}

\vspace{0.05cm}
\texttt{linguagem natural}

\vspace{0.05cm}
{\color{orange}\large\downarrow}

\vspace{0.05cm}
\textbf{delegação de trabalho intelectual}
\end{center}

\vfill{\tiny\color{muted}A ideia dialoga com cognição distribuída e descarregamento cognitivo: Risko e Gilbert (2016), DOI 10.1016/j.tics.2016.07.002.}
\end{frame}

\begin{frame}{O ganho é real - e é por isso que o risco importa}
\protect\phantomsection\label{o-ganho-uxe9-real---e-uxe9-por-isso-que-o-risco-importa}
\begin{columns}[T,onlytextwidth]
\column{0.48\textwidth}
\Large\textbf{A IA pode}

\vspace{0.25cm}
\begin{itemize}
\item explicar;
\item calcular;
\item cruzar dados;
\item gerar gráficos;
\item explorar alternativas.
\end{itemize}

\column{0.48\textwidth}
\Large\textbf{O pesquisador precisa}

\vspace{0.25cm}
\begin{itemize}
\item definir o objeto;
\item escolher a medida;
\item assumir as hipóteses;
\item interpretar o contraste;
\item defender a conclusão.
\end{itemize}
\end{columns}

\vfill{\tiny\color{muted}A dependência de ferramentas não é automaticamente perda de autonomia; o critério é preservar a autodireção intelectual. Carter (2020), DOI 10.1007/s11229-017-1549-y.}
\end{frame}

\begin{frame}{Onde a abstração começa a cobrar juros}
\protect\phantomsection\label{onde-a-abstrauxe7uxe3o-comeuxe7a-a-cobrar-juros}
\begin{center}
\vspace{0.45cm}
\Large A especificação pede “estime a dimensão”.

\vspace{0.45cm}
\large A IA pode completar silenciosamente:

\vspace{0.35cm}
\begin{tabular}{c@{\quad}c@{\quad}c}
  \color{orange}\textbf{o que medir} & \color{orange}\textbf{quem conta} & \color{orange}\textbf{o que concluir} \\
  objeto & população/denominador & estimando/interpretação
\end{tabular}

\vspace{0.6cm}
\textbf{O código pode estar correto. A pergunta pode ter mudado.}
\end{center}

\vfill{\tiny\color{muted}A “fronteira” da IA é irregular: a assistência melhora algumas tarefas e piora outras, e o usuário nem sempre sabe previamente quais são quais. Dell’Acqua et al. (2026), DOI 10.1287/orsc.2025.21838.}
\end{frame}

\begin{frame}{Minha proposta: dívida epistêmica}
\protect\phantomsection\label{minha-proposta-duxedvida-epistuxeamica}
\begin{center}
\vspace{0.6cm}
{\color{orange}\Huge\textbf{Dívida epistêmica}}

\vspace{0.45cm}
\Large decisões substantivas que ficaram incorporadas ao artefato,

\vspace{0.15cm}
\Large mas deixaram de estar compreendidas e assumidas pelo autor.

\vspace{0.65cm}
\normalsize\textit{Formulação autoral, apoiada em literaturas adjacentes.}
\end{center}
\end{frame}

\begin{frame}{O caso: pontos ideais e as âncoras}
\protect\phantomsection\label{o-caso-pontos-ideais-e-as-uxe2ncoras}
\begin{center}
\vspace{0.25cm}
\large\textbf{“As âncoras foram definidas teoricamente antes dos resultados - ou porque produziam uma dimensão conveniente?”}

\vspace{0.7cm}
\begin{tabular}{c@{\qquad}c@{\qquad}c}
  \color{orange}\textbf{proveniência} & \color{orange}\textbf{mensuração} & \color{orange}\textbf{defesa} \\
  Por que escolhi? & O que a dimensão mede? & Consigo explicar?
\end{tabular}
\end{center}

\vfill{\tiny\color{muted}A literatura de pontos ideais trata identificação e interpretação como problemas distintos: Poole e Rosenthal (1985); Clinton, Jackman e Rivers (2004); Morucci et al. (2025), DOI 10.1017/S000305542400039X.}
\end{frame}

\begin{frame}{O que exatamente foi perdido?}
\protect\phantomsection\label{o-que-exatamente-foi-perdido}
\begin{columns}[T,onlytextwidth]
\column{0.31\textwidth}
\centering
{\color{orange}\Huge 1}

\vspace{0.12cm}
\Large\textbf{Proveniência}

\vspace{0.18cm}
\normalsize Não reconstruía a razão da escolha.

\column{0.31\textwidth}
\centering
{\color{orange}\Huge 2}

\vspace{0.12cm}
\Large\textbf{Construto}

\vspace{0.18cm}
\normalsize Não estava claro o que a dimensão representava.

\column{0.31\textwidth}
\centering
{\color{orange}\Huge 3}

\vspace{0.12cm}
\Large\textbf{Autoria}

\vspace{0.18cm}
\normalsize O resultado existia; a defesa não.
\end{columns}

\vfill{\tiny\color{muted}Um ajuste latente não garante significado substantivo: Morucci et al. (2025); Adcock e Collier (2001), DOI 10.1017/S0003055401003100.}
\end{frame}

\begin{frame}{Reprodução ajuda - mas não resolve tudo}
\protect\phantomsection\label{reproduuxe7uxe3o-ajuda---mas-nuxe3o-resolve-tudo}
\begin{center}
\vspace{0.35cm}
\Large
\textbf{Código e logs} \quad $\neq$ \quad \textbf{validade} \quad $\neq$ \quad \textbf{autoria}

\vspace{0.65cm}
\begin{columns}[T,onlytextwidth]
\column{0.31\textwidth}
\centering\textbf{Reproduzir}

\vspace{0.18cm}
\normalsize o que foi executado

\column{0.31\textwidth}
\centering\textbf{Validar}

\vspace{0.18cm}
\normalsize o que foi medido

\column{0.31\textwidth}
\centering\textbf{Defender}

\vspace{0.18cm}
\normalsize por que essa escolha é minha
\end{columns}

\vspace{0.55cm}
\large Salvaguardas são necessárias. Elas não substituem deliberação.
\end{center}

\vfill{\tiny\color{muted}Diretrizes recentes recomendam registrar ferramenta, versão, datas, prompts e procedimentos para permitir reprodução: Flanagin et al. (2024), DOI 10.1001/jama.2024.3471; Abdurahman et al. (2025), DOI 10.1177/25152459251325174.}
\end{frame}

\begin{frame}{A defesa metodológica sem cola}
\protect\phantomsection\label{a-defesa-metodoluxf3gica-sem-cola}
\begin{center}
\vspace{0.2cm}
\large
\textbf{IA implementa} $\rightarrow$ \textbf{IA questiona} $\rightarrow$ \textbf{pesquisador investiga} $\rightarrow$ \textbf{pesquisador decide}

\vspace{0.35cm}
\small
\begin{enumerate}
\setlength{\itemsep}{2pt}
\item Peça um Q\&A adversarial.
\item Não peça imediatamente as respostas.
\item Transforme cada pergunta sem resposta em pendência.
\item Reconstrua a proveniência e compare alternativas.
\item Delibere, registre e assuma a decisão.
\end{enumerate}
\end{center}

\vfill{\tiny\color{muted}A lógica é próxima das intervenções de “forçamento cognitivo”, que reduziram sobredelegação em decisões assistidas por IA: Buçinca, Malaya e Gajos (2021), DOI 10.1145/3449287.}
\end{frame}

\begin{frame}{O meu limiar provisório}
\protect\phantomsection\label{o-meu-limiar-provisuxf3rio}
\begin{center}
\vspace{0.5cm}
{\color{teal}\Huge\textbf{Delegue a implementação,}}

\vspace{0.2cm}
{\color{orange}\Huge\textbf{não a autoria da cadeia inferencial.}}

\vspace{0.75cm}
\large Pause e delibere quando uma alternativa razoável puder mudar:

\vspace{0.2cm}
\normalsize objeto \quad | \quad denominador \quad | \quad operacionalização \quad | \quad estimando \quad | \quad interpretação
\end{center}
\end{frame}

\begin{frame}{O que muda na prática?}
\protect\phantomsection\label{o-que-muda-na-pruxe1tica}
\begin{center}
\vspace{0.35cm}
\begin{columns}[T,onlytextwidth]
\column{0.46\textwidth}
\Large\textbf{Antes de aceitar}

\vspace{0.2cm}
\normalsize “A IA sugeriu.”

\vspace{0.18cm}
\normalsize “O código roda.”

\vspace{0.18cm}
\normalsize “O resultado parece plausível.”

\column{0.46\textwidth}
\Large\textbf{Depois de aceitar}

\vspace{0.2cm}
\normalsize “Eu comparei alternativas e escolhi.”

\vspace{0.18cm}
\normalsize “O objeto e o estimando estão claros.”

\vspace{0.18cm}
\normalsize “Consigo responder à crítica.”
\end{columns}
\end{center}

\vfill{\tiny\color{muted}A literatura recente descreve riscos de ilusão de entendimento, opacidade de proveniência e confiança mal calibrada em pesquisa assistida por IA: Messeri e Crockett (2024); Gautam et al. (2026).}
\end{frame}

\begin{frame}{Fechamento}
\protect\phantomsection\label{fechamento}
\begin{center}
\vspace{0.75cm}
{\color{teal}\Huge\textbf{Seu paper não precisa ser feito sem IA.}}

\vspace{0.35cm}
{\color{orange}\Huge\textbf{Precisa continuar sendo defendível por você.}}

\vspace{0.8cm}
\large A pergunta final não é “a IA fez?”.

\vspace{0.15cm}
\large É: \textbf{“qual decisão eu assumi - e por quê?”}
\end{center}
\end{frame}

\begin{frame}{Para aprofundar}
\protect\phantomsection\label{para-aprofundar}
\small

\begin{itemize}
\item Risko e Gilbert (2016). “Cognitive Offloading.” \textit{Trends in Cognitive Sciences}. DOI: 10.1016/j.tics.2016.07.002.
\item Carter (2020). “Intellectual Autonomy, Epistemic Dependence and Cognitive Enhancement.” \textit{Synthese}. DOI: 10.1007/s11229-017-1549-y.
\item Messeri e Crockett (2024). “Artificial Intelligence and Illusions of Understanding in Scientific Research.” \textit{Nature}. DOI: 10.1038/s41586-024-07146-0.
\item Morucci et al. (2025). “Measurement That Matches Theory.” \textit{APSR}. DOI: 10.1017/S000305542400039X.
\item Abdurahman et al. (2025). “A Primer for Evaluating LLMs in Social-Science Research.” DOI: 10.1177/25152459251325174.
\end{itemize}

\vfill
\centering

\textcolor{muted}{A revisão exploratória completa e as referências estão no arquivo de apoio do projeto.}
\end{frame}

\end{document}
